% Options for packages loaded elsewhere
\PassOptionsToPackage{unicode}{hyperref}
\PassOptionsToPackage{hyphens}{url}
\documentclass[
]{article}
\usepackage{xcolor}
\usepackage{amsmath,amssymb}
\setcounter{secnumdepth}{-\maxdimen} % remove section numbering
\usepackage{iftex}
\ifPDFTeX
  \usepackage[T1]{fontenc}
  \usepackage[utf8]{inputenc}
  \usepackage{textcomp} % provide euro and other symbols
\else % if luatex or xetex
  \usepackage{unicode-math} % this also loads fontspec
  \defaultfontfeatures{Scale=MatchLowercase}
  \defaultfontfeatures[\rmfamily]{Ligatures=TeX,Scale=1}
\fi
\usepackage{lmodern}
\ifPDFTeX\else
  % xetex/luatex font selection
\fi
% Use upquote if available, for straight quotes in verbatim environments
\IfFileExists{upquote.sty}{\usepackage{upquote}}{}
\IfFileExists{microtype.sty}{% use microtype if available
  \usepackage[]{microtype}
  \UseMicrotypeSet[protrusion]{basicmath} % disable protrusion for tt fonts
}{}
\makeatletter
\@ifundefined{KOMAClassName}{% if non-KOMA class
  \IfFileExists{parskip.sty}{%
    \usepackage{parskip}
  }{% else
    \setlength{\parindent}{0pt}
    \setlength{\parskip}{6pt plus 2pt minus 1pt}}
}{% if KOMA class
  \KOMAoptions{parskip=half}}
\makeatother
\usepackage{longtable,booktabs,array}
\newcounter{none} % for unnumbered tables
\usepackage{calc} % for calculating minipage widths
% Correct order of tables after \paragraph or \subparagraph
\usepackage{etoolbox}
\makeatletter
\patchcmd\longtable{\par}{\if@noskipsec\mbox{}\fi\par}{}{}
\makeatother
% Allow footnotes in longtable head/foot
\IfFileExists{footnotehyper.sty}{\usepackage{footnotehyper}}{\usepackage{footnote}}
\makesavenoteenv{longtable}
\usepackage{graphicx}
\makeatletter
\newsavebox\pandoc@box
\newcommand*\pandocbounded[1]{% scales image to fit in text height/width
  \sbox\pandoc@box{#1}%
  \Gscale@div\@tempa{\textheight}{\dimexpr\ht\pandoc@box+\dp\pandoc@box\relax}%
  \Gscale@div\@tempb{\linewidth}{\wd\pandoc@box}%
  \ifdim\@tempb\p@<\@tempa\p@\let\@tempa\@tempb\fi% select the smaller of both
  \ifdim\@tempa\p@<\p@\scalebox{\@tempa}{\usebox\pandoc@box}%
  \else\usebox{\pandoc@box}%
  \fi%
}
% Set default figure placement to htbp
\def\fps@figure{htbp}
\makeatother
\ifLuaTeX
  \usepackage{luacolor}
  \usepackage[soul]{lua-ul}
\else
  \usepackage{soul}
\fi
\setlength{\emergencystretch}{3em} % prevent overfull lines
\providecommand{\tightlist}{%
  \setlength{\itemsep}{0pt}\setlength{\parskip}{0pt}}
\usepackage{bookmark}
\IfFileExists{xurl.sty}{\usepackage{xurl}}{} % add URL line breaks if available
\urlstyle{same}
\hypersetup{
  hidelinks,
  pdfcreator={LaTeX via pandoc}}

\author{}
\date{}

\begin{document}

\begin{quote}
\includegraphics[width=0.72917in,height=0.72917in,alt={Graphic Era logo}]{./media/image1.jpg}
\end{quote}

\section{\texorpdfstring{PROJECT AND TEAM INFORMATION
}{PROJECT AND TEAM INFORMATION }}\label{project-and-team-information}

Project Title

\begin{quote}
RESQNET: Smart Urban Emergency \& Transportation Network Simulator
\end{quote}

{\def\LTcaptype{none} % do not increment counter
\begin{longtable}[]{@{}
  >{\raggedright\arraybackslash}p{(\linewidth - 2\tabcolsep) * \real{0.5000}}
  >{\raggedright\arraybackslash}p{(\linewidth - 2\tabcolsep) * \real{0.5000}}@{}}
\toprule\noalign{}
\begin{minipage}[b]{\linewidth}\raggedright
\emph{Team Name:} \emph{Team \#}
\end{minipage} & \begin{minipage}[b]{\linewidth}\raggedright
\emph{RESQNET}

\emph{DSCPP-III-2026-T070}
\end{minipage} \\
\midrule\noalign{}
\endhead
\bottomrule\noalign{}
\endlastfoot
\emph{\textbf{Team member 1 (Team Lead)}}

\emph{(Last Name, name: student ID: email, picture):} &
\begin{minipage}[t]{\linewidth}\raggedright
\emph{Chhavi -- 2510012506}

\emph{\ul{2510012506@geu.ac.in}}

\begin{quote}
\includegraphics[width=0.625in,height=0.78125in,alt={photo}]{./media/image2.jpg}
\end{quote}
\end{minipage} \\
\emph{\textbf{Team member 2}}

\emph{(Last Name, name: student ID: email, picture):} &
\begin{minipage}[t]{\linewidth}\raggedright
\emph{Tapashya -- 2510011649}

\emph{\ul{2510011649@geu.ac.in}}

\begin{quote}
\includegraphics[width=0.625in,height=0.875in,alt={photo}]{./media/image3.jpg}
\end{quote}
\end{minipage} \\
\emph{\textbf{Team member 3}}

\emph{(Last Name, name: student ID: email, picture):} &
\begin{minipage}[t]{\linewidth}\raggedright
\emph{Chahal, Arpit -- 2510012000}

\emph{\ul{2510012000@geu.ac.in}}

\begin{quote}
\includegraphics[width=0.625in,height=0.70833in,alt={photo}]{./media/image4.jpg}
\end{quote}
\end{minipage} \\
\end{longtable}
}

\section{\texorpdfstring{PHASE-II PROJECT PROGRESS DESCRIPTION (35 pts)
}{PHASE-II PROJECT PROGRESS DESCRIPTION (35 pts) }}\label{phase-ii-project-progress-description-35-pts}

Project Abstract (2 pts)

\begin{quote}
RESQNET is a C++17 console-based simulator that models how a city
responds to several emergencies at the same time. The city is
represented as a weighted graph in which junctions, hospitals and key
locations are nodes and roads are edges carrying distance and
travel-time weights. The simulator prioritizes incoming incidents with a
priority queue, selects suitable emergency vehicles from a shared
resource pool, and routes them to the incident using Dijkstra's
algorithm.

The overall goal is one integrated model that shows the connected
decisions of emergency response: incident prioritization, resource
allocation, traffic-aware routing, dispatch, rerouting around road
failures, hospital selection and critical-junction analysis. The project
focuses on algorithmic decision-making using Data Structures and OOP
concepts, not on live GPS or hardware.
\end{quote}

Updated Project Approach and Architecture (2 pts)

\begin{quote}
The system follows the layered architecture from our proposal: an input
layer, a core system layer and a data layer, with a user interface (UI)
layer for simulation controls and output display. The core flow is
Disaster Event → Incident Analysis → Priority Engine → Resource
Allocation → Routing Engine → Dispatch \& Response → Hospital Selection
→ Network Analysis → Logging.

Currently implemented: the weighted city graph, the Event Manager and
incident handling, the Priority Engine (priority queue), the Resource
Allocation engine, the Routing Engine (Dijkstra) and the Dispatch \&
Response step.

Data structures and algorithms: weighted graph, priority queue, vectors,
queues and hash maps; Dijkstra for routing, with BFS/DFS and
articulation points planned for network analysis. OOP is used for
Vehicle, Disaster and Hospital classes with inheritance, polymorphism
and encapsulation.

Libraries and tools: C++17 with the Standard Template Library only (no
external libraries); VS Code / Code::Blocks as IDEs; Git and GitHub for
version control. The application runs as a single process, so modules
communicate through function calls and shared in-memory data and no
network protocol is used.
\end{quote}

Tasks Completed (7 pts)

{\def\LTcaptype{none} % do not increment counter
\begin{longtable}[]{@{}
  >{\raggedright\arraybackslash}p{(\linewidth - 2\tabcolsep) * \real{0.6313}}
  >{\raggedright\arraybackslash}p{(\linewidth - 2\tabcolsep) * \real{0.3687}}@{}}
\toprule\noalign{}
\begin{minipage}[b]{\linewidth}\raggedright
\textbf{Task Completed}
\end{minipage} & \begin{minipage}[b]{\linewidth}\raggedright
\textbf{Team Member}
\end{minipage} \\
\midrule\noalign{}
\endhead
\bottomrule\noalign{}
\endlastfoot
Problem statement finalized; disaster and vehicle types identified;
system architecture and class structure designed & All members (Chhavi
as Team Lead) \\
City road network implemented in C++ as a weighted graph (junctions,
roads, distance and travel-time weights) & Chhavi \\
Dijkstra's shortest-path routing implemented & Chhavi \\
Disaster/incident handling and emergency priority queue implemented &
Chahal, Arpit \\
Resource allocation (selecting a suitable available vehicle for an
incident) and OOP classes & Tapashya \\
Dispatch of selected resources to the incident & Tapashya \\
User interface (UI) for running and viewing the simulation & Chhavi \\
\end{longtable}
}

Challenges/Roadblocks (7 pts)

{\def\LTcaptype{none} % do not increment counter
\begin{longtable}[]{@{}
  >{\raggedright\arraybackslash}p{(\linewidth - 2\tabcolsep) * \real{0.0479}}
  >{\raggedright\arraybackslash}p{(\linewidth - 2\tabcolsep) * \real{0.9521}}@{}}
\toprule\noalign{}
\begin{minipage}[b]{\linewidth}\raggedright
\begin{quote}
•
\end{quote}
\end{minipage} & \begin{minipage}[b]{\linewidth}\raggedright
Keeping routes correct when edge weights change. Travel cost must react
to traffic and road failures, so the plan is to keep all weights inside
one graph object and recompute Dijkstra whenever a weight changes.
\end{minipage} \\
\midrule\noalign{}
\endhead
\bottomrule\noalign{}
\endlastfoot
\begin{minipage}[t]{\linewidth}\raggedright
\begin{quote}
•
\end{quote}
\end{minipage} & Integrating modules written by different members
(priority queue, allocation, routing, dispatch). We agreed on common
class interfaces and merge work in small steps using Git branches. \\
\begin{minipage}[t]{\linewidth}\raggedright
\begin{quote}
•
\end{quote}
\end{minipage} & Several incidents competing for the same vehicles.
Allocation uses incident priority together with vehicle availability and
distance, and will be tested with many simultaneous incidents. \\
\end{longtable}
}

Tasks Pending (7 pts)

{\def\LTcaptype{none} % do not increment counter
\begin{longtable}[]{@{}
  >{\raggedright\arraybackslash}p{(\linewidth - 2\tabcolsep) * \real{0.6313}}
  >{\raggedright\arraybackslash}p{(\linewidth - 2\tabcolsep) * \real{0.3687}}@{}}
\toprule\noalign{}
\begin{minipage}[b]{\linewidth}\raggedright
\textbf{Task Pending}
\end{minipage} & \begin{minipage}[b]{\linewidth}\raggedright
\textbf{Team Member (to complete the task)}
\end{minipage} \\
\midrule\noalign{}
\endhead
\bottomrule\noalign{}
\endlastfoot
Hospital selection by capability, capacity and travel time & Tapashya \\
Dynamic rerouting on road or junction failure & Chhavi \\
Critical junction detection using articulation points & Chhavi \\
Traffic congestion model and time-based simulation clock & Chahal,
Arpit \\
File logging (fstream) of events and results & Chahal, Arpit \\
Scenario testing, comparison of results, documentation and final demo &
All members \\
\end{longtable}
}

Project Outcome/Deliverables (2 pts)

{\def\LTcaptype{none} % do not increment counter
\begin{longtable}[]{@{}
  >{\raggedright\arraybackslash}p{(\linewidth - 2\tabcolsep) * \real{0.0479}}
  >{\raggedright\arraybackslash}p{(\linewidth - 2\tabcolsep) * \real{0.9521}}@{}}
\toprule\noalign{}
\begin{minipage}[b]{\linewidth}\raggedright
\begin{quote}
•
\end{quote}
\end{minipage} & \begin{minipage}[b]{\linewidth}\raggedright
A complete C++17 console simulator with modules for emergency response
management, smart traffic management and dynamic navigation / network
analysis.
\end{minipage} \\
\midrule\noalign{}
\endhead
\bottomrule\noalign{}
\endlastfoot
\begin{minipage}[t]{\linewidth}\raggedright
\begin{quote}
•
\end{quote}
\end{minipage} & The city modeled as a weighted graph of roads,
junctions, hospitals and emergency points, with a multi-vehicle resource
pool. \\
\begin{minipage}[t]{\linewidth}\raggedright
\begin{quote}
•
\end{quote}
\end{minipage} & Source code, project report/documentation with system
design and flow diagrams, test results and the final project
presentation. \\
\end{longtable}
}

Progress Overview (2 pts)

\begin{quote}
Of the 13 milestones in our proposal, milestones 1 (simulation rules), 2
(base city map) and 4 (emergency resources and allocation) are complete,
and milestone 3 (incident lifecycle) is in progress: creation,
prioritization and dispatch work, while arrival and closing remain. The
core pipeline Disaster → Priority Queue → Resource Allocation → Dijkstra
Routing → Dispatch is working.

The Phase II targets from our roadmap (graph, priority queue, Dijkstra
routing, resource allocation) are met, and dispatch and the UI were also
completed. Milestones 5 to 13 (traffic and congestion, intersection
priority, rerouting, hospital selection, network analysis, full
simulation, scenario comparison and final testing) are pending and
planned for the next phase.
\end{quote}

Codebase Information (2 pts)

\begin{quote}
Repository link: https://github.com/mchhavi929-ship-it/Resqnet Branch:
main
\end{quote}

Testing and Validation Status (2 pts)

{\def\LTcaptype{none} % do not increment counter
\begin{longtable}[]{@{}
  >{\raggedright\arraybackslash}p{(\linewidth - 4\tabcolsep) * \real{0.2806}}
  >{\raggedright\arraybackslash}p{(\linewidth - 4\tabcolsep) * \real{0.1503}}
  >{\raggedright\arraybackslash}p{(\linewidth - 4\tabcolsep) * \real{0.5691}}@{}}
\toprule\noalign{}
\begin{minipage}[b]{\linewidth}\raggedright
\textbf{Test Type}
\end{minipage} & \begin{minipage}[b]{\linewidth}\raggedright
\begin{quote}
\textbf{Status}

\textbf{(Pass/Fail)}
\end{quote}
\end{minipage} & \begin{minipage}[b]{\linewidth}\raggedright
\begin{quote}
\textbf{Notes}
\end{quote}
\end{minipage} \\
\midrule\noalign{}
\endhead
\bottomrule\noalign{}
\endlastfoot
Graph construction & \begin{minipage}[t]{\linewidth}\raggedright
\begin{quote}
Pass
\end{quote}
\end{minipage} & \begin{minipage}[t]{\linewidth}\raggedright
\begin{quote}
City graph built with junctions and roads; node count, edges and
distance/travel-time weights checked against the input data.
\end{quote}
\end{minipage} \\
Dijkstra shortest route & \begin{minipage}[t]{\linewidth}\raggedright
\begin{quote}
Pass
\end{quote}
\end{minipage} & \begin{minipage}[t]{\linewidth}\raggedright
\begin{quote}
Routes between junction pairs compared with manually calculated shortest
paths on a small sample graph.
\end{quote}
\end{minipage} \\
Priority queue ordering & \begin{minipage}[t]{\linewidth}\raggedright
\begin{quote}
Pass
\end{quote}
\end{minipage} & \begin{minipage}[t]{\linewidth}\raggedright
\begin{quote}
Incidents of different severity inserted in random order; the
highest-priority incident is always removed first.
\end{quote}
\end{minipage} \\
Multiple incident handling & \begin{minipage}[t]{\linewidth}\raggedright
\begin{quote}
Pass
\end{quote}
\end{minipage} & \begin{minipage}[t]{\linewidth}\raggedright
\begin{quote}
Several simultaneous incidents of different disaster types created and
processed in priority order.
\end{quote}
\end{minipage} \\
Resource allocation & \begin{minipage}[t]{\linewidth}\raggedright
\begin{quote}
Pass
\end{quote}
\end{minipage} & \begin{minipage}[t]{\linewidth}\raggedright
\begin{quote}
Selected vehicle matches the disaster type, is available and is the best
choice by distance.
\end{quote}
\end{minipage} \\
Dispatch & \begin{minipage}[t]{\linewidth}\raggedright
\begin{quote}
Pass
\end{quote}
\end{minipage} & \begin{minipage}[t]{\linewidth}\raggedright
\begin{quote}
Dispatched vehicle is marked busy and is not assigned to another
incident until it is free.
\end{quote}
\end{minipage} \\
Integrated flow (end to end) &
\begin{minipage}[t]{\linewidth}\raggedright
\begin{quote}
Pass
\end{quote}
\end{minipage} & \begin{minipage}[t]{\linewidth}\raggedright
\begin{quote}
Disaster → priority queue → allocation → routing → dispatch run with
sample city data.
\end{quote}
\end{minipage} \\
Dynamic rerouting, hospital selection &
\begin{minipage}[t]{\linewidth}\raggedright
\begin{quote}
Pending
\end{quote}
\end{minipage} & \begin{minipage}[t]{\linewidth}\raggedright
\begin{quote}
Not tested yet; these features are still to be implemented.
\end{quote}
\end{minipage} \\
\end{longtable}
}

Deliverables Progress (2 pts)

{\def\LTcaptype{none} % do not increment counter
\begin{longtable}[]{@{}
  >{\raggedright\arraybackslash}p{(\linewidth - 2\tabcolsep) * \real{0.5210}}
  >{\raggedright\arraybackslash}p{(\linewidth - 2\tabcolsep) * \real{0.4790}}@{}}
\toprule\noalign{}
\begin{minipage}[b]{\linewidth}\raggedright
\textbf{Deliverable}
\end{minipage} & \begin{minipage}[b]{\linewidth}\raggedright
\begin{quote}
\textbf{Status}
\end{quote}
\end{minipage} \\
\midrule\noalign{}
\endhead
\bottomrule\noalign{}
\endlastfoot
C++ simulator: graph, priority queue, Dijkstra routing, resource
allocation, dispatch, UI & \begin{minipage}[t]{\linewidth}\raggedright
\begin{quote}
Completed
\end{quote}
\end{minipage} \\
C++ simulator: hospital selection, rerouting, traffic model,
articulation points, simulation clock, logging &
\begin{minipage}[t]{\linewidth}\raggedright
\begin{quote}
Pending
\end{quote}
\end{minipage} \\
System design and flow diagrams &
\begin{minipage}[t]{\linewidth}\raggedright
\begin{quote}
Completed (architecture diagram); flow diagrams in progress
\end{quote}
\end{minipage} \\
Project report / documentation &
\begin{minipage}[t]{\linewidth}\raggedright
\begin{quote}
In progress
\end{quote}
\end{minipage} \\
Test results and scenario comparison &
\begin{minipage}[t]{\linewidth}\raggedright
\begin{quote}
Pending
\end{quote}
\end{minipage} \\
Project presentation & \begin{minipage}[t]{\linewidth}\raggedright
\begin{quote}
Phase I completed; final presentation pending
\end{quote}
\end{minipage} \\
\end{longtable}
}

\end{document}
